\documentclass[12pt,a4paper]{article}
\usepackage{fontspec}
\setmainfont{Times New Roman}
\usepackage[top=2.0cm,bottom=2.0cm,left=2.5cm,right=2.0cm]{geometry}
\usepackage{enumitem}
\usepackage{tabularx}
\usepackage{xltabular}
\usepackage{booktabs}
\usepackage{array}
\newcolumntype{Y}{>{\centering\arraybackslash}X}
\newcolumntype{C}[1]{>{\centering\arraybackslash}p{#1}}
\newcolumntype{L}[1]{>{\raggedright\arraybackslash}p{#1}}
\newcolumntype{R}[1]{>{\raggedleft\arraybackslash}p{#1}}
\pagestyle{plain}
\linespread{1.15}
\setlength{\parskip}{4pt plus 1pt minus 1pt}
\sloppy

\begin{document}
% --- HEADER 2 CỘT CHUẨN NGHỊ ĐỊNH 30/2020/NĐ-CP ---
\noindent
\begin{minipage}[t]{0.46\textwidth}
    \begin{center}
        \fontsize{11pt}{13pt}\selectfont
        CÔNG AN THÀNH PHỐ HÀ NỘI\\
        \textbf{PHÒNG CẢNH SÁT HÌNH SỰ (PC02)}\\[-0.2em]
        \rule{3.2cm}{0.6pt}\\[0.25cm]
        \fontsize{11pt}{13pt}\selectfont Số: \textbf{00/HS}
    \end{center}
\end{minipage}
\hfill
\begin{minipage}[t]{0.52\textwidth}
    \begin{center}
        \fontsize{11pt}{13pt}\selectfont
        \textbf{CỘNG HÒA XÃ HỘI CHỦ NGHĨA VIỆT NAM}\\
        \textbf{Độc lập - Tự do - Hạnh phúc}\\[-0.2em]
        \rule{3.6cm}{0.6pt}\\[0.25cm]
        \textit{Hà Nội, ngày 25 tháng 07 năm 2016}
    \end{center}
\end{minipage}

\vspace{0.5cm}

\begin{center}
    \textbf{\Large BIÊN BẢN KHÁM XÉT}\\[0.15cm]
    \textit{(Về việc khám xét khẩn cấp chỗ ở)}
\end{center}

\vspace{0.3cm}

\textit{Căn cứ Lệnh khám xét khẩn cấp số 45/LKX-PC02 ngày 25/07/2016 của Thủ trưởng Cơ quan Cảnh sát điều tra, Công an TP. Hà Nội đối với chỗ ở của đối tượng Trần Thị Hà.}


Vào hồi 14 giờ 30 phút, ngày 26 tháng 07 năm 2016.


Tại: Phòng trọ tại số 08, ngõ 12, đường Bờ Kè, Phường Phân khu Cảng, Quận Sông Hồng, TP. Hà Nội.


\textbf{Thành phần tiến hành khám xét gồm:}

\begin{itemize}[leftmargin=*, itemsep=2pt, topsep=2pt, parsep=0pt]
  \item Điều tra viên chủ trì: \textbf{Lê Minh}. Chức vụ: Đại úy, Điều tra viên Phòng PC02 — Công an TP. Hà Nội.
  \item Cán bộ ghi biên bản: \textbf{Nguyễn Văn Hoàng}. Chức vụ: Trung úy, Điều tra viên Phòng PC02.

\end{itemize}
\textbf{Thành phần tham gia và chứng kiến khám xét gồm:}


\textbf{1. Người có chỗ ở bị khám xét:}


Đối tượng TRẦN THỊ HÀ.    Sinh năm: 1990.     Nơi ở: Phòng trọ nêu trên.


\textbf{2. Đại diện chính quyền địa phương:}


Ông Nguyễn Văn Bình      Chức vụ: Công an Phường Phân khu Cảng, Quận Sông Hồng.


\textbf{3. Người chứng kiến:}

\begin{itemize}[leftmargin=*, itemsep=2pt, topsep=2pt, parsep=0pt]
  \item Ông: Đỗ Văn Vinh     Sinh năm: 1968.     Trú tại: Số 10 ngõ 12, đường Bờ Kè, Phường Phân khu Cảng, Quận Sông Hồng.
  \item Bà: Lê Thị Thảo      Sinh năm: 1975.     Trú tại: Số 14 ngõ 12, đường Bờ Kè, Phường Phân khu Cảng, Quận Sông Hồng.


\end{itemize}

\vspace{0.3cm}\noindent\textbf{\large KẾT QUẢ KHÁM XẾT VÀ THU GIỮ ĐỒ VẬT, TÀI LIỆU}\\[0.15cm]


Tiến hành khám xét toàn bộ diện tích phòng trọ và tư trang của đối tượng Trần Thị Hà. Cơ quan điều tra phát hiện và thu giữ các đồ vật, tài liệu sau:


\begin{enumerate}[label=\arabic*., leftmargin=*, itemsep=2pt, topsep=2pt, parsep=0pt]
  \item \textbf{01 lọn tóc nam giới:} Dài khoảng 4,2cm, chân tóc dính vết chất màu nâu khô. Lọn tóc được chứa trong 01 túi ni lông nhỏ trong suốt (túi zip), cất giấu bên trong ví tiền của đối tượng Trần Thị Hà.
  \item \textbf{01 chiếc áo gió nam có mũ trùm đầu:} Màu xám đen, vải dù. Áo nằm trong tủ quần áo của đối phượng. Quan sát bề mặt áo phát hiện dính bẩn phấn hoa.
  \item \textbf{01 chiếc kéo bấm kim loại nhỏ:} Dạng kéo cắt chỉ thêu len, dài khoảng 10cm, để trên bàn làm việc cạnh cuộn len. Tại phần lưỡi kéo phát hiện có vệt chất màu nâu khô và sợi tóc bám dính.
  \item \textbf{01 lá bùa:} Có kích thước, chất liệu tương đồng với lá bùa đã tìm thấy tại nhà riêng nạn nhân.

\end{enumerate}
Ngoài các đồ vật nêu trên, Cơ quan điều tra không thu giữ thêm tài liệu, đồ vật nào khác. Toàn bộ các đồ vật thu giữ đã được kiểm tra, niêm phong tại chỗ trước sự chứng kiến của đối tượng Trần Thị Hà, đại diện chính quyền địa phương và những người chứng kiến.


Việc khám xét kết thúc vào hồi 15 giờ 30 phút cùng ngày. Biên bản này đã được đọc lại cho tất cả những người có tên nêu trên nghe, công nhận chính xác và cùng ký tên xác nhận dưới đây.



\textbf{NGƯỜI CÓ CHỖ Ở BỊ KHÁM XẾT}

\textit{(Ký, ghi rõ họ tên)}


\textit{(Đã ký)}

\textbf{Trần Thị Hà}


\textbf{ĐẠI DIỆN CHÍNH QUYỀN ĐỊA PHƯƠNG}

\textit{(Ký, ghi rõ họ tên)}


\textit{(Đã ký)}

\textbf{Nguyễn Văn Bình}


\textbf{CÁC NGƯỜI CHỨNG KIẾN}

\textit{(Ký, ghi rõ họ tên)}


\textit{(Đã ký)}

\vspace{0.8cm}
\noindent
\begin{minipage}[t]{0.48\textwidth}
    \begin{center}
        \textbf{Đỗ Văn Vinh}\\[0.1cm]
        \textit{(Ký, ghi rõ họ tên)}\\[1.4cm]
        \textbf{Lê Thị Thảo}
    \end{center}
\end{minipage}
\hfill
\begin{minipage}[t]{0.48\textwidth}
    \begin{center}
        \textbf{CÁN BỘ GHI BIÊN BẢN}\\[0.1cm]
        \textit{(Ký, ghi rõ họ tên)}\\[1.4cm]
        \textbf{Nguyễn Văn Hoàng}
    \end{center}
\end{minipage}
\vspace{0.6cm}
\noindent
\hfill
\begin{minipage}[t]{0.50\textwidth}
    \begin{center}
        \textbf{ĐIỀU TRA VIÊN CHỦ TRÌ}\\[0.1cm]
        \textit{(Ký, ghi rõ họ tên)}\\[1.4cm]
        \textbf{Lê Minh}
    \end{center}
\end{minipage}
\end{document}
